\documentclass{article}
\usepackage[utf8]{inputenc}
\usepackage[french]{babel}
\usepackage{amsmath}
\usepackage{amsfonts}
\usepackage{amssymb}
\usepackage{geometry}
\geometry{a4paper, margin=1in}

% Raccourcis
\newcommand{\pref}[2]{s_{#1:#2}}                     % préfixe d'indices s_{a:b}
\newcommand{\w}[2]{\omega\big(#1 \mid #2\big)}       % poids SNIS
\newcommand{\E}{\mathbb{E}}
\newcommand{\Prob}{\mathbb{P}}

\begin{document}

\title{Dérivation SNIS multi-tour :\\
pourquoi l'échantillonnage ancestral absorbe les poids}
\author{Implémentation \texttt{src/train/train\_grpo\_snis.py}}
\date{9 juillet 2026}
\maketitle

\section{Cadre et notations}

On fixe un prompt $x$ et le \textbf{pool} des $G$ trajectoires réelles générées par
GRPO. Toutes les quantités ci-dessous sont conditionnelles à ce pool : les tours
$z^k_j$ sont des objets \emph{fixés} (des séquences de tokens déjà générées), et la
seule source d'aléa considérée ici est le choix des \textbf{indices}
$s_k \in \{1,\dots,G\}$. On note un préfixe d'indices
$\pref{1}{k-1} = (s_1, \dots, s_{k-1})$, et on définit le poids SNIS du tour $k$ :
\begin{equation}
\w{s_k}{\pref{1}{k-1}}
\;=\;
\frac{\pi_\theta\big(z^{k}_{s_k} \,\big|\, x,\; z^1_{s_1}, \ldots, z^{k-1}_{s_{k-1}}\big)}
     {\displaystyle\sum_{l=1}^{G} \pi_\theta\big(z^{k}_{l} \,\big|\, x,\; z^1_{s_1}, \ldots, z^{k-1}_{s_{k-1}}\big)}
\label{eq:poids}
\end{equation}
ainsi que la fonction dont on veut la moyenne (le terme le plus interne de la
formule forward) :
\begin{equation}
f(s_1, \ldots, s_N)
\;=\;
\nabla_\theta \log \pi_\theta\big((z^k_{s_k})_{k=1}^{N} \,\big|\, x\big)
\cdot r\big(x,\, z^N_{s_N}\big).
\label{eq:f}
\end{equation}

La formule forward de la note s'écrit alors :
\begin{equation}
\hat{\nabla}_\theta V
\;=\;
\frac{1}{G} \sum_{s_1=1}^{G}\;
\underbrace{\sum_{s_2=1}^{G} \w{s_2}{s_1}
\sum_{s_3=1}^{G} \w{s_3}{\pref{1}{2}} \cdots
\sum_{s_N=1}^{G} \w{s_N}{\pref{1}{N-1}}\;
f(s_1, \ldots, s_N)}_{=:\;S(s_1)}.
\label{eq:forward}
\end{equation}
Tout le travail porte sur le bloc interne $S(s_1)$, qui contient $G^{N-1}$ termes.

\section{Étape 1 --- À chaque niveau, les poids somment à 1}

Pour \emph{n'importe quel} préfixe $\pref{1}{k-1}$ fixé :
\begin{equation}
\sum_{s_k=1}^{G} \w{s_k}{\pref{1}{k-1}}
= \frac{\sum_{s_k=1}^{G} \pi_\theta\big(z^k_{s_k} \mid x, z^1_{s_1},\ldots\big)}
       {\sum_{l=1}^{G} \pi_\theta\big(z^k_{l} \mid x, z^1_{s_1},\ldots\big)}
= 1.
\label{eq:normalisation}
\end{equation}
C'est trivial --- le dénominateur \emph{est} la somme des numérateurs --- mais c'est
\textbf{la} propriété structurelle du self-normalized IS : sans le dénominateur, les
$\pi(z^k_l \mid \cdot)$ ne sommeraient à rien de particulier (ce sont $G$
évaluations de densité, pas une partition de l'espace). Avec lui, et comme
$\pi_\theta > 0$, on a pour chaque préfixe fixé :
\begin{center}
$\w{\cdot}{\pref{1}{k-1}}$ est une loi de probabilité sur $\{1, \ldots, G\}$.
\end{center}

\section{Étape 2 --- Le produit des poids est une loi jointe sur les combinaisons}

Pour $s_1$ fixé, on définit :
\begin{equation}
P\big(s_2, \ldots, s_N \mid s_1\big)
\;=\;
\prod_{k=2}^{N} \w{s_k}{\pref{1}{k-1}}.
\label{eq:loiP}
\end{equation}

\paragraph{Affirmation.} $P(\cdot \mid s_1)$ est une loi de probabilité sur
l'ensemble des $G^{N-1}$ combinaisons $(s_2,\ldots,s_N) \in \{1,\ldots,G\}^{N-1}$.

\paragraph{Preuve.} Positivité : produit de termes strictement positifs.
Normalisation : on somme, et on fait sortir les sommes une par une \emph{en
commençant par la plus interne}. L'argument qui l'autorise : $\w{s_N}{\pref{1}{N-1}}$
est le \emph{seul} facteur du produit qui dépend de $s_N$, donc :
\begin{equation}
\sum_{s_2, \ldots, s_N} \prod_{k=2}^{N} \w{s_k}{\pref{1}{k-1}}
= \sum_{s_2, \ldots, s_{N-1}} \left[ \prod_{k=2}^{N-1} \w{s_k}{\pref{1}{k-1}} \right]
  \underbrace{\sum_{s_N} \w{s_N}{\pref{1}{N-1}}}_{=\,1 \;\;\text{par \eqref{eq:normalisation}}}.
\end{equation}
Le point subtil : la somme interne vaut $1$ \textbf{quelle que soit la valeur du
préfixe} $\pref{1}{N-1}$ (l'équation \eqref{eq:normalisation} est valable pour tout
préfixe) --- c'est ce qui permet de la remplacer par $1$ \emph{à l'intérieur} de la
somme externe. On répète l'argument : la somme sur $s_{N-1}$ vaut alors $1$ pour
tout $\pref{1}{N-2}$, etc. Le télescope descend jusqu'à :
\begin{equation}
\sum_{s_2} \w{s_2}{s_1} = 1
\qquad\Longrightarrow\qquad
\sum_{s_2, \ldots, s_N} P\big(s_2,\ldots,s_N \mid s_1\big) = 1. \qquad \blacksquare
\end{equation}

\section{Étape 3 --- La somme imbriquée EST une espérance sous $P$}

Le bloc $S(s_1)$ de la formule \eqref{eq:forward} est exactement la moyenne de $f$
pesée par $P$. On distribue les poids vers l'intérieur : chaque $\w{s_k}{\cdot}$ est
une constante vis-à-vis des sommes sur $s_{k+1},\ldots,s_N$ situées à sa droite,
donc il entre dedans :
\begin{align}
S(s_1)
&= \sum_{s_2} \w{s_2}{s_1} \sum_{s_3} \w{s_3}{\pref{1}{2}} \cdots
   \sum_{s_N} \w{s_N}{\pref{1}{N-1}}\, f(\pref{1}{N}) \\
&= \sum_{s_2, \ldots, s_N} \left[ \prod_{k=2}^{N} \w{s_k}{\pref{1}{k-1}} \right]
   f(s_1, s_2, \ldots, s_N)
\end{align}
soit, par définition de $P$ \eqref{eq:loiP} puis de l'espérance d'une variable
discrète :
\begin{equation}
\boxed{\;
S(s_1)
= \sum_{s_2, \ldots, s_N} P\big(\pref{2}{N} \mid s_1\big)\, f(s_1, \pref{2}{N})
= \E_{\,\pref{2}{N} \sim P(\cdot \mid s_1)}\big[\, f(s_1, \pref{2}{N}) \,\big].
\;}
\label{eq:esperance}
\end{equation}
Rien d'approximatif jusqu'ici : c'est une \textbf{réécriture exacte} de la formule
\eqref{eq:forward}. Elle ne devient utile qu'avec les étapes 4 et 5.

\section{Étape 4 --- Estimateur Monte-Carlo de cette espérance}

Principe général : si $P$ est une loi sur un ensemble fini et
$S^{(1)}, \ldots, S^{(M)} \overset{iid}{\sim} P$, alors
\begin{equation}
\widehat{E}_M = \frac{1}{M} \sum_{m=1}^{M} f\big(S^{(m)}\big),
\qquad
\E\big[\widehat{E}_M\big] = \sum_{s} P(s)\, f(s) = \E_P[f]
\end{equation}
--- estimateur \textbf{sans biais quel que soit $M$}, y compris $M=1$. On applique
ceci à chaque $S(s_1)$ avec $M=1$ (un seul tirage par valeur de $s_1$ : contrainte
des $G$ sorties attendues par TRL) :
\begin{equation}
\widehat{\nabla_\theta V}
= \frac{1}{G} \sum_{s_1=1}^{G} f\big(s_1,\, S^{(s_1)}\big),
\qquad S^{(s_1)} \sim P(\cdot \mid s_1),
\label{eq:estimateur}
\end{equation}
et par linéarité de l'espérance :
\begin{equation}
\E\Big[\widehat{\nabla_\theta V} \;\Big|\; \text{pool}\Big]
= \frac{1}{G} \sum_{s_1=1}^{G} \E_{P(\cdot \mid s_1)}[f]
\overset{\eqref{eq:esperance}}{=} \frac{1}{G} \sum_{s_1=1}^{G} S(s_1)
\overset{\eqref{eq:forward}}{=} \hat{\nabla}_\theta V
\quad \text{(la formule, exactement).}
\end{equation}

\paragraph{Hiérarchie des deux niveaux d'estimation.} Notre tirage est
\textbf{sans biais pour la formule} \eqref{eq:forward}, conditionnellement au pool ;
la formule SNIS est elle-même un estimateur \textbf{biaisé en $O(1/G)$} du vrai
gradient (le dénominateur de \eqref{eq:poids} est une \emph{estimation} de la
constante de normalisation, pas sa valeur exacte --- biais classique du SNIS). Le
Monte-Carlo n'ajoute que de la variance, pas de biais.

\section{Étape 5 --- L'échantillonnage ancestral tire exactement de $P$}

Comment obtenir $S^{(s_1)} \sim P(\cdot \mid s_1)$ sans énumérer les $G^{N-1}$
combinaisons ? Par la \textbf{règle de chaînage} : on tire séquentiellement
\begin{equation}
s_2 \sim \w{\cdot}{s_1}, \qquad
s_3 \sim \w{\cdot}{\pref{1}{2}}, \qquad \ldots, \qquad
s_N \sim \w{\cdot}{\pref{1}{N-1}},
\end{equation}
chaque tirage utilisant le préfixe \emph{déjà réalisé} aux étapes précédentes.
La loi jointe du vecteur obtenu est, par construction du tirage séquentiel :
\begin{equation}
\Prob\big(s_2 = a_2, \ldots, s_N = a_N\big)
= \prod_{k=2}^{N} \Prob\big(s_k = a_k \mid s_2 = a_2, \ldots, s_{k-1} = a_{k-1}\big)
= \prod_{k=2}^{N} \w{a_k}{s_1, \pref{2}{k-1} = a_{2:k-1}}
= P\big(a_{2:N} \mid s_1\big).
\end{equation}
C'est un tirage \textbf{exact} (pas approché) parce que $P$ a été \emph{définie}
sous forme factorisée : ses conditionnelles sont précisément les $\omega$.
Coût : $G$ scorings par niveau, soit $O(N \cdot G)$ scorings par combo au lieu de
$G^N$ termes.

\section{Étape 6 --- « Absorption des poids » : les deux estimateurs côte à côte}

Deux façons d'obtenir la même espérance $\E_P[f]$ :
\begin{equation}
\text{(a) somme exhaustive : } \sum_{s} P(s)\, f(s)
\qquad\qquad
\text{(b) Monte-Carlo : } f(S), \;\; S \sim P.
\end{equation}
En (a), les poids apparaissent \textbf{explicitement comme coefficients} --- c'est
la formule \eqref{eq:forward}, $G^{N-1}$ termes. En (b), les poids n'apparaissent
\textbf{que dans le tireur} ; la quantité utilisée est $f(S)$ avec coefficient $1$.
L'égalité $\E[f(S)] = \sum_s P(s) f(s)$ dit que (b) estime (a) sans qu'aucun poids
ne survive dans l'expression finale.

Traduction dans le pipeline d'entraînement : la quantité que TRL calcule pour chaque
combo retourné est $f(\text{combo}) = \nabla_\theta \log \pi_\theta(\text{combo})
\cdot r$, avec un poids uniforme dans le batch --- c'est exactement la forme (b).
Voilà pourquoi la loss GRPO n'a pas été modifiée : tout le SNIS vit dans le tirage
pondéré, et plus aucune trace de $\omega$ ne sort de la fonction de recombinaison.

\section{Hypothèses utilisées par cette dérivation (et où elle craque)}

\begin{enumerate}
  \item \textbf{$N$ identique pour toutes les trajectoires.} La dérivation suppose
  que chaque somme $\sum_{s_k}$ porte sur les mêmes $G$ candidats. En pratique les
  épisodes ont des longueurs différentes : le code restreint le support aux
  trajectoires possédant un tour $k$, et termine un combo quand il consomme le
  dernier tour de sa source. L'équation \eqref{eq:normalisation} reste vraie sur le
  support restreint, donc le télescope tient --- mais la \emph{loi} échantillonnée
  n'est plus exactement celle de la formule idéalisée.
  \item \textbf{$f$ ne dépend des indices que via les tours recousus.} Vrai pour
  $\nabla \log \pi$ ; pour $r$, on utilise l'hypothèse de la note ($r$ fonction de
  $x$ et du dernier tour seulement) --- c'est l'approximation « reward hérité de la
  source du dernier tour ».
  \item \textbf{$r$ remplacé par l'avantage GRPO.} TRL centre et réduit les rewards
  par groupe de $G$. Le baseline $b$ ajoute un terme
  $-b \cdot \frac{1}{G}\sum_m \nabla \log \pi(\text{combo}_m)$ dont l'espérance
  sous l'échantillonnage SNIS n'est qu'\emph{approximativement} nulle (elle l'est
  exactement en on-policy pur). Petit biais supplémentaire.
\end{enumerate}

\section*{Correspondance formule $\leftrightarrow$ code}

\begin{center}
\begin{tabular}{l l}
\hline
\textbf{Objet mathématique} & \textbf{Code (\texttt{src/train/train\_grpo\_snis.py})} \\
\hline
Poids $\omega$ \eqref{eq:poids} (normalisation) & \texttt{\_snis\_weights}, l.~304--312 \\
Numérateurs $\pi(z^k_l \mid \text{préfixe})$ & jobs l.~373--377, \texttt{score\_fn} l.~380 \\
$\log \pi(\text{séquence}) = \sum_t \log \pi(\text{token}_t)$ & l.~392 \\
Support de $\sum_{s_k}$ (candidats du tour $k$) & l.~370 \\
Énumération de $\frac{1}{G}\sum_{s_1}$ (combo $m \Rightarrow s_1 = m$) & l.~333--346 \\
Un niveau $\sum_{s_k}$ = une itération & boucle \texttt{while}, l.~348--409 \\
Tirage $s_k \sim \w{\cdot}{\text{préfixe}}$ (étape 5) & \texttt{rng.choices}, l.~399--400 \\
Extension du préfixe (fixer $s_k$) & \texttt{\_append\_turn}, l.~404--405 \\
$\nabla \log \pi \times$ avantage (forme (b)) & TRL \texttt{grpo\_trainer.py}, l.~2479--2498 \\
Masquage des observations dans $\log \pi$ & \texttt{env\_mask} $\to$ TRL l.~1762, l.~2327 \\
Avantage par groupe de $G$ & TRL l.~2146--2168 \\
\hline
\end{tabular}
\end{center}

\end{document}
